\chapter{Methods}\label{ch:methods}

This chapter describes the design of the hybrid reconstruction pipeline:
the property graph schema that stores the static artefacts, the separate
trace-collection pipeline, the agent skill architecture that orchestrates
extraction and generation, the pattern-based instrumentation tool used for
trace collection, and the scenario-scoped querying that produces
scenario-specific interaction models.

\section{Pipeline Overview}\label{sec:pipeline-overview}

The pipeline takes a \Cpp\ source tree and, optionally, a set of
timestamped execution traces, and produces \sysml sequence diagrams
scoped to a user-specified scenario. It runs as six stages
(\cref{fig:pipeline-overview}), each implemented as a composable agent
skill invocable by natural-language query and backed by a shared \neo
property graph. Five stages follow a fixed, deterministic rule set given
their inputs; the sixth (Stage~5, Agentic Synthesis) is the one skill
whose own instructions are reasoning guidance rather than a strict rule
set which resists a fixed lookup.

\begin{figure}[htb]
\centering
\begin{tikzpicture}[
  node distance = 8mm and 14mm,
  stage/.style = {rectangle, draw, rounded corners, thick,
                  minimum width=2.6cm, minimum height=1.05cm,
                  align=center, font=\small},
  dyn/.style   = {stage, fill=blue!15},
  stat/.style  = {stage, fill=gray!12},
  agentic/.style = {stage, fill=orange!20},
  arr/.style   = {-{Stealth[length=2.2mm]}, thick}
]
  \node[stat] (extract)  {1. Static\\Extraction};
  \node[stat, right=of extract] (query) {2. Code Graph\\Query};
  \node[stat, right=of query]   (matrix) {3. Dependency\\Matrix Extraction};
  \node[dyn,  below=of matrix]  (trace)  {4. Trace Instrumentation\\\& Runtime Tracing};
  \node[agentic, left=of trace] (synth)  {5. Agentic\\Synthesis};
  \node[stat, left=of synth]    (sysml)  {6. SysML~v2\\Generation};

  \draw[arr] (extract) -- (query);
  \draw[arr] (query)   -- (matrix);
  \draw[arr] (matrix)  -- (trace);
  \draw[arr] (trace)   -- (synth);
  \draw[arr] (synth)   -- (sysml);

  \begin{scope}[on background layer]
    \node[draw, dashed, fit=(extract)(query)(matrix),
          inner sep=6mm, label={[font=\small]above:Static graph construction \& scoping}] {};
  \end{scope}
\end{tikzpicture}
\caption{Pipeline overview. Numbered stages correspond to the enumeration
  below. Stage~4 (blue) is the only stage that reads dynamic (runtime)
  information directly; the static graph and the trace-derived table
  are kept separate up to this point and combined only here
  (\cref{sec:combine-scope,sec:agentic-synthesis}).}
\label{fig:pipeline-overview}
\end{figure}

\begin{enumerate}
  \item \textbf{Static extraction} (\cref{sec:static-impl}). The \Cpp\
    source is analysed by Renaissance, which populates the \neo graph
    with 11 node labels and 11 relationship types (\cref{sec:pg-schema});
    the CPSCore case-study graph contains 67,104 nodes.

  \item \textbf{Code graph query} (\texttt{text-to-cypher}, built on
    \texttt{get-neo4j-schema}; \texttt{cypher-to-output}). A
    natural-language scenario description is translated into a
    schema-aware Cypher query by \texttt{text-to-cypher} and executed by
    \texttt{cypher-to-output} against the live graph, retaining only
    interactions whose caller and callee belong to the scenario's
    declared components (\cref{sec:scoping-impl}).
    \texttt{get-neo4j-schema} is queried first so query generation
    targets the live schema rather than hardcoded node labels, keeping
    the pipeline codebase-agnostic.

  \item \textbf{Dependency matrix extraction}
    (\texttt{get-interface-dependency-table},
    \texttt{get-sequence-dependency-table}). The scenario-scoped subgraph
    is flattened into an ordered CSV table, one row per \texttt{CppCalls}
    edge, giving the static ground truth against which the runtime trace
    is later checked.

  \item \textbf{Trace instrumentation and runtime tracing}
    (\texttt{add-runtime-instrumentation}). Pattern-based code
    transformation injects trace probes at the call sites named in the
    dependency matrix; the instrumented binary is rebuilt and run, and
    the emitted \texttt{[TRACE]} lines are normalised into the same table
    schema.

  \item \textbf{Agentic synthesis} (\texttt{agentic-synthesis},
    \cref{sec:agentic-synthesis}). An \ac{LLM} is given the static
    dependency matrix and the runtime trace as two separately-labelled,
    un-merged evidence sources and settles the final interaction set by
    judgement rather than a fixed rule, writing the result to
    \texttt{elements.json} and \texttt{sequence.json}.

  \item \textbf{SysML v2 generation} (\texttt{draw-sysml-sequence-model},
    \texttt{draw-sysml-structural-model}). \texttt{elements.json} and
    \texttt{sequence.json} are projected into \sysml textual notation.
\end{enumerate}

Two further skills, \texttt{draw-sequential-mermaid-general} and
\texttt{draw-structural-mermaid}, offer a Mermaid-syntax visualisation of
the same reconstruction; they read the Stage~3 dependency tables directly
and are not part of this six-stage numbering, since they bypass agentic
synthesis rather than consuming its output.

Each skill is stateless and read-only with respect to the graph, except
the instrumentation skill, which modifies source files. Every skill
accepts a natural-language prompt and returns a structured artefact (CSV,
Cypher query, or model file), which allows the chain to be reused on a
new codebase without code changes. The agentic synthesis skill is the
one exception, per \cref{sec:pipeline-overview}: its judgement calls
(\cref{sec:agentic-synthesis}) resist a deterministic procedure.

\section{Property Graph Schema}\label{sec:pg-schema}

\subsection{Design Decision}\label{sec:pg-decision}

A \ac{PG} in \neo was chosen for the static knowledge store over a
relational database, a document store, and an in-memory graph library.
The pipeline traverses function$\to$file$\to$folder paths,
function$\to$function call chains, and class$\to$function containment
within the same query; a typed, directed property graph represents each
as a first-class edge, where a relational schema would need join tables.
The static graph can also be re-derived or extended without disturbing
the separate trace pipeline, since static edges are typed and
distinguishable from anything added later, and Cypher's pattern-matching
syntax maps directly onto the \texttt{text-to-cypher} skill
(\cref{sec:skill-design}): a natural-language scenario query becomes a
graph pattern with little translation. An in-memory graph library would
satisfy traversal alone, and a document store neither traversal nor
declarative querying.

\subsection{Node Types}\label{sec:node-types}

\Cref{tab:node-types} lists the 11 node labels emitted by Renaissance.

\begin{table}[htb]
\centering
\caption{Property graph node labels relevant to reconstruction. Each C++
  node (the \texttt{Cpp*} labels) carries \texttt{name}, \texttt{symbol}, \texttt{type}, and
  \texttt{lineNumber} properties}
\label{tab:node-types}
\small
\begin{tabular}{@{} l p{9cm} @{}}
\toprule
\textbf{Label} & \textbf{Semantics} \\
\midrule
\texttt{CppDeclaration}           & Class, struct, enum, or typedef declaration \\
\texttt{CppForwardDeclaration}    & Forward declaration of a class or struct \\
\texttt{CppFunctionDeclaration}   & Function prototype (header) \\
\texttt{CppFunctionDefinition}    & Function body (implementation) \\
\texttt{CppMacroDefinition}       & Preprocessor macro definition \\
\texttt{HeaderFile}               & \texttt{.h}/\texttt{.hpp} file \\
\texttt{SourceFile}               & \texttt{.cpp}/\texttt{.c} file \\
\texttt{OtherFile}                & Non-C++ file tracked in the project \\
\texttt{Folder}                   & Directory in the source tree \\
\texttt{Symbol}                   & Demangled symbol name (linked to declarations via \texttt{Name} edges) \\
\texttt{ProjectCOrCpp}            & CMake compilation unit \\
\bottomrule
\end{tabular}
\end{table}

Renaissance additionally emits four build-system labels
(\texttt{DynamicLinkLibrary}, \texttt{Executable},
\texttt{ProjectPropertySheet}, \texttt{Solution}) omitted from the table:
they carry no interaction-reconstruction information and the pipeline
never queries them. The graph holds static artefacts only.

\subsection{Relationship Types}\label{sec:edge-types}

\Cref{tab:edge-types} lists the 11 relationship types populated by
Renaissance.

\begin{table}[htb]
\centering
\caption{Property graph relationship types, all populated by Renaissance
  from static analysis.}
\label{tab:edge-types}
\small
\begin{tabular}{@{} l l l p{5.5cm} @{}}
\toprule
\textbf{Type} & \textbf{Source} & \textbf{Target} & \textbf{Semantics} \\
\midrule
\texttt{CppCalls}       & FunctionDef & FunctionDecl/Def & Direct call site \\
\texttt{CppContains}    & Declaration & Declaration/FunctionDecl/Def & Lexical containment \\
\texttt{CppImplements}  & FunctionDef & FunctionDecl/Def & Declaration-definition link \\
\texttt{CppInherits}    & Declaration & Declaration/ForwardDecl & Class inheritance \\
\texttt{CppAlias}       & Declaration & Declaration & Type alias \\
\texttt{CppUses}        & File        & Declaration/ForwardDecl/Macro & Symbol usage in file \\
\texttt{Include}        & File        & File & Preprocessor include \\
\texttt{Name}           & Cpp* node   & Symbol & Name binding \\
\texttt{ParentFolder}   & File/Folder & Folder & Directory hierarchy \\
\texttt{ProjectCompilesNested} & ProjectCOrCpp & File & Compilation membership \\
\texttt{Source}         & Cpp* node   & File & Source-location back-reference \\
\bottomrule
\end{tabular}
\end{table}

The static \texttt{CppCalls} edges and the trace-derived sequence table
are kept as two independent edge sets and compared only at evaluation
time (\cref{sec:combine-scope}): a \texttt{CppCalls} edge says a call
\emph{may} happen, since it is read from source code listing every call
that could occur, while a trace record says a call \emph{did} happen,
recorded from a real run. The static edges therefore tend to over-count
and the trace tends to under-count.

\section{Agent Skill Architecture}\label{sec:skill-design}

The pipeline is a set of twelve composable agent skills rather than a
monolithic script. Each is defined by a declarative specification (name,
description, input/output contracts, execution instructions) and
invocable via natural-language prompt. This was chosen over a monolithic
implementation for three reasons:

\begin{itemize}
  \item \textbf{Codebase agnosticism.} Skills query the graph schema
    dynamically (\texttt{get-neo4j-schema}) rather than hardcoding node
    labels, so switching target codebases requires only re-running
    Renaissance.
  \item \textbf{Composability.} Each skill produces an intermediate
    artefact (a CSV, a Cypher query, a model file) that an engineer can
    inspect, edit, or override before it feeds the next skill.
  \item \textbf{Scenario flexibility.} The scope of a reconstruction can
    be expressed as a natural-language query (``show me the interaction
    between Synchronization and Aggregation when a runner timeout
    occurs'') rather than a hardcoded configuration, with
    \texttt{text-to-cypher} translating it into the appropriate Cypher
    pattern at invocation time.
\end{itemize}

\Cref{tab:skills} summarises all twelve skills, their inputs, outputs,
and role in the pipeline; only the agentic synthesis skill departs from
a fixed rule set (\cref{sec:pipeline-overview,sec:agentic-synthesis}).

\begin{table}[htb]
\centering
\caption{Agent skills composing the reconstruction pipeline.}
\label{tab:skills}
\small
{\renewcommand{\arraystretch}{1.6}
\begin{tabularx}{\textwidth}{@{} >{\ttfamily}p{0.25\textwidth} p{0.20\textwidth} p{0.25\textwidth} >{\raggedright\arraybackslash}X @{}}
\toprule
\textbf{Skill} & \textbf{Input} & \textbf{Output} & \textbf{Role} \\
\midrule
\texttt{get-neo4j-schema}
  & Neo4j connection
  & \texttt{graph\_schema.txt}
  & Schema introspection \\
\texttt{text-to-cypher}
  & Schema + NL query
  & Cypher query
  & Query generation \\
\texttt{cypher-to-output}
  & Cypher query
  & Human-readable answer
  & Query execution + explanation \\
\texttt{get-interface-dependency-table}
  & Neo4j graph
  & Interface CSV
  & Structural extraction \\
\texttt{get-sequence-dependency-table}
  & Neo4j graph
  & Sequence CSV
  & Sequence (\texttt{CppCalls}) extraction \\
\texttt{add-runtime-instrumentation}
  & Source files + patterns
  & Instrumented source + Trace CSV
  & Probe injection via csp\_matcher, rebuild, run, and trace collection \\
\texttt{agentic-synthesis}
  & Static table + trace (CSV or JSON)
  & \texttt{elements.json} + \texttt{sequence.json}
  & Evidence adjudication (reasoning guidance, not a fixed rule set) \\
\texttt{draw-sysml-sequence-model}
  & Sequence CSV or \texttt{sequence.json}
  & \texttt{.sysml} file
  & Sequence diagram generation \\
\texttt{draw-sysml-structural-model}
  & Interface CSV or \texttt{elements.json}
  & \texttt{.sysml} file
  & Structural diagram generation \\
\texttt{draw-sequential-mermaid-general}
  & Sequence CSV
  & \texttt{.mmd} file
  & Mermaid sequence visualisation \\
\texttt{draw-structural-mermaid}
  & Interface CSV
  & \texttt{.mmd} file
  & Mermaid dependency visualisation \\
\texttt{get-cypher}
  & NL query
  & Cypher (no execution)
  & Standalone query authoring \\
\bottomrule
\end{tabularx}}
\end{table}

\section{Pattern-Based Instrumentation (csp\_matcher)}\label{sec:csp-matcher}

Automated insertion of trace probes into an existing \Cpp\ codebase is a
prerequisite for runtime trace collection: manual instrumentation is
error-prone and does not scale, and a preliminary experiment with an LLM
alone proved brittle, since the model could not reliably identify every
call site of interest and tended to over-instrument. The
\texttt{csp\_matcher} tool, developed alongside Dr.~Rosilde Corvino as
part of this thesis, addresses this with AST-level pattern matching and
code transformation.

\subsection{Design}

\texttt{csp\_matcher} is a Clang LibTooling-based tool that parses a
user-provided pattern (near-normal \Cpp\ with hole variables) into a
Clang \ac{AST}, translates it into ASTMatcher DSL expressions, runs the
generated matchers against target source files, and optionally rewrites
matched code using a template that back-references captured holes.

Patterns use two kinds of hole variables:
\begin{description}
  \item[\texttt{\$name}] Matches any single expression or identifier.
  \item[\texttt{\$\$name}] Matches a statement list or parameter list.
\end{description}

Two context variables are automatically injected for every match:
\texttt{\$callerFunc} (the enclosing function name) and
\texttt{\$callerClass} (the enclosing class name, empty at file scope).

Example: wrap every \texttt{getAll} call with a trace statement:

\begin{minted}{text}
find:    $agg.getAll<$T>()
replace: (std::fprintf(stderr, "[TRACE] %s %s::%s\n",
           cspTraceTimestamp().c_str(),
           "$callerClass", "$callerFunc"),
          $agg.getAll<$T>())
\end{minted}

This is preferable to manual macro insertion in three respects: the
transformation is declarative, so a rules file states \emph{what} to
match and replace rather than \emph{how} to walk the AST; it is
reproducible, since re-running the same rules file produces the same
edits; and it is reversible, via a corresponding \texttt{remove:} rule
(\cref{sec:reversibility}). Filter functions, compiled at runtime from
plain \Cpp, additionally allow selective matching, \eg excluding
test-internal call sites.

\subsection{Integration with the Pipeline}

The \texttt{add-runtime-instrumentation} skill invokes
\texttt{csp\_matcher} with a rules file (\texttt{add\_tracing.txt}) and a
filter definitions file (\texttt{tracing\_filters.cpp}), following a
two-pronged approach. First, call-site wrapping: target call expressions
such as \texttt{Aggregator::getAll}, \texttt{CPSLogger::flush}, and
\texttt{EnumMap::convert} are wrapped with a comma-operator trace
emission. Second, CPSLOG macro patching: the \texttt{CPSLOG(level)}
macro in \texttt{CPSLogger.h} is patched to emit a \texttt{[TRACE]} line
on every logging call, covering all \texttt{RAIILogStream::stream} call
sites regardless of translation unit. The macro patch is applied only
\emph{after} call-site wrapping is complete, since \texttt{csp\_matcher}
parses target files using a pattern preamble that includes
\texttt{CPSLogger.h}, which must be in its clean, unpatched state when
the matcher runs.

Once the probes are in place, the same skill rebuilds and runs the
target (the CPSCore test suite, or a scenario-specific harness) and
captures the emitted \texttt{[TRACE]} lines from \texttt{stderr}. This
log is normalised into the same schema as the static dependency matrix
(\cref{sec:pipeline-overview}) so the two can later be compared
edge-for-edge: it does not pass through
\texttt{get-sequence-dependency-table}, which extracts static
\texttt{CppCalls} edges from the property graph rather than reading
trace output; the normalised trace log instead feeds the evidence
combination and agentic synthesis stages directly
(\cref{sec:combine-scope,sec:agentic-synthesis}).

\subsection{Reversibility via Filter Functions}\label{sec:reversibility}

Instrumentation removal is handled by a single generic rule:

\begin{minted}[fontsize=\footnotesize]{text}
find:    ($any, $orig)
filter:  is_trace_remove
replace: $orig
\end{minted}

This matches every comma expression in the codebase, but the
\texttt{is\_trace\_remove} filter rejects any match whose left operand
does not contain the \texttt{[TRACE]} marker string:

\begin{minted}[fontsize=\footnotesize]{cpp}
bool is_trace_remove(int n, const char* const* ns,
                     const char* const* vs) {
    const char* text = get_hole(n, ns, vs, "any");
    return text != NULL && strstr(text, "[TRACE]") != NULL;
}
\end{minted}

Only trace-injected comma expressions are stripped; pre-existing comma
expressions are left untouched, so applying
\texttt{remove\_tracing.txt} restores the source tree to its
pre-instrumentation state byte-for-byte. \texttt{csp\_matcher} was
validated by an automated test suite of 130 cases across seven suites
(\cref{sec:appendix-tracing-rules}).

\section{Macro Resolution for Publish-Subscribe Systems}\label{sec:macro-resolution}

The CPSCore IDC layer (\cref{sec:idc}) implements publish-subscribe via
\texttt{boost::signals2} callbacks registered through
\texttt{subscribeOnPacket}. A static call graph sees an edge from the
subscribing component into the IDC layer, but not from the publisher to
the subscriber's handler, since the two sides are wired by a shared
runtime channel identifier rather than a direct call.

Macro resolution recovers this topology by matching publisher and
subscriber declarations on that identifier: every
\texttt{subscribeOnPacket(id, slot)} and \texttt{sendPacket(id, payload)}
call site is located in the static graph, and pairs sharing the same
\texttt{id} are linked by a synthetic \texttt{PubSubRoute} edge. This
lets the sequence dependency table contain cross-component interactions
invisible to naive call-graph extraction, and generalises to any
publish-subscribe system keyed by a shared identifier (topic name,
channel ID, signal slot).

This only resolves \texttt{subscribeOnPacket}/\texttt{sendPacket} pairs
sharing an \texttt{id}. A pair wired directly through a
\texttt{boost::signals2::signal} and \texttt{connect()}, outside the IDC
abstraction, has no shared identifier to match and is observable only in
a runtime trace (S4, \cref{sec:h1-c4c3}).

\section{Combining Evidence and Scenario Scoping}\label{sec:combine-scope}

\subsection{Combining the Two Evidence Sources}

The static and dynamic evidence live in two separate artefacts: the
static \texttt{CppCalls} edges in the property graph, and the
trace-derived sequence table (\cref{sec:csp-matcher}). Each is projected
to a set of directed, component-level interaction edges keyed by
\texttt{(ClientComponent, ServerComponent, EventName)}, and the two edge
sets are combined only when the reconstructed models are evaluated.

The static edge set is obtained with a plain Cypher query:

\begin{minted}{cypher}
MATCH (caller:CppFunctionDefinition)-[:Source]->(f)-[:ParentFolder]->(folder)
MATCH (caller)-[:CppCalls]->(callee)
  WHERE last(split(folder.name, '/')) = $component
RETURN caller, callee
\end{minted}

The dynamic edge set is read from the real runtime trace, normalised
into the same schema. The union of the two, tagged by which source each
interaction came from, is what condition C4 of the H1 experiment operates
over directly (\cref{ch:experiments}), a set operation over two separate
artefacts computed at evaluation time. Condition C5, the full
pipeline, takes this same tagged union as the input to the agentic
synthesis stage (\cref{sec:agentic-synthesis}), which prunes it further
before SysML generation rather than passing it through unchanged.

\subsection{Scenario Scoping by Component Constraint}\label{sec:scoping-impl}

Scenario scoping is a plain, scope-constrained Cypher query rather than
the program-dependence-graph slice described in \cref{sec:scenario-scoping}.
A scenario's scope is a conjunction of a
\emph{primary component} constraint (only edges whose caller belongs to
it are retained), a \emph{target component set} (only edges whose callee
belongs to a declared participant is retained), and optionally a
\emph{client function filter} restricting to specific entry-point
functions. For scenario S1, concretely:

\begin{minted}{text}
ClientComponent == Synchronization
ServerComponent in {Aggregation, Logging}
\end{minted}

This is what \texttt{text-to-cypher} generates from a scenario
description such as ``Using agent skills, show me Synchronization
calling into Aggregation and Logging''. The explicit \emph{Using agent
skills} prefix anchors the model to the skill invocation, preventing it
from hallucinating an ad-hoc codebase inspection instead of generating a
Cypher query.

The same constraint applies to both evidence sources: a \texttt{WHERE}
clause over the static graph, and an equivalent row filter over the
trace-derived table (both carry \texttt{ClientComponent} and
\texttt{ServerComponent} metadata). Since the scope is defined over
components, it can be computed \emph{before} instrumentation and used to
decide where to place probes (\cref{sec:guided-vs-full}), rather than
applied only as a post-hoc filter.

\subsection{Formal Description}

Let $E_s$ be the static interaction edge set and $E_t$ the trace-derived
edge set, with $\text{src}(e)$ and $\text{tgt}(e)$ an edge's caller and
callee. Let $\sigma = (\text{primary}, \text{targets}, \text{functions})$
be the scenario scope. The scenario-scoped, combined edge set is:
\[
  E_\sigma
  = \{ e \in E_s \cup E_t \mid \text{src}(e) \in \sigma.\text{primary}
       \land \text{tgt}(e) \in \sigma.\text{targets} \}
\]
that is, $E_s \cup E_t$ restricted to edges originating from the primary
component and terminating in a declared target component. This is a
single-hop edge filter, which suffices for CPSCore, where cross-component
calls are shallow (typically one hop through the IDC layer); multi-hop
scoping is a straightforward extension via Cypher's variable-length path
syntax.

\section{Agentic Synthesis}\label{sec:agentic-synthesis}

The tagged union $E_s \cup E_t$ (\cref{sec:combine-scope}) is not yet the
final interaction set: a static edge with no trace corroboration is
either a genuine over-approximation (a call the source code permits but
this scenario never takes) or a real interaction that test coverage
simply missed, and a dynamic edge outside the static scope may be a real
indirect-dispatch interaction the call graph could not see
(\cref{sec:macro-resolution}) rather than noise. Telling these apart
requires judgement a set operation cannot make on its own, so the union
is passed to an \ac{LLM} (GPT-5, invoked via Azure OpenAI) as the final
pipeline stage before diagram generation.

\subsection{Mechanism}

The \ac{LLM} receives the static and dynamic evidence as two separately
labelled, un-merged sources (it is never shown a pre-merged edge list)
and reasons over three judgement categories:

\begin{description}
  \item[Corroboration.] An interaction present in both sources is kept
    without further scrutiny.
  \item[Suppression.] A static-only interaction is kept if it is
    plausibly real but untested by this scenario's runtime coverage, and
    dropped if it reads as a structural over-approximation.
  \item[Indirect dispatch.] A dynamic-only interaction is kept if it
    reflects a real call path the static graph could not resolve, such as
    a \texttt{boost::signals2} connection made outside the \ac{IDC}
    abstraction (S4, \cref{sec:h1-c4c3}).
\end{description}

The output is written to two files: \texttt{elements.json}, an array of
deduplicated \texttt{(source, target, interaction)} triples forming the
final structural element set, and \texttt{sequence.json}, the same
interactions in execution order, each step optionally carrying a
free-text \texttt{note} field. The \texttt{note} field is populated only
when a step's evidence came from a single source rather than being
corroborated both ways, so an empty \texttt{note} marks a trivial
judgement call and a populated one marks a step where the agent actually
exercised discretion.

\subsection{Where the Judgement Is Trivial versus Non-Trivial}

For scenario S1, every static edge in scope is corroborated by the trace
and vice versa, so the agent's task reduces to relabelling an
already-agreeing union: all steps carry an empty \texttt{note}. The
harder cases are where Stage~5 has to do something non-trivial with the
freedom it is given, for instance S4's suppressed fault-path log entry
and the \texttt{boost::signals2} indirect-dispatch edge
(\cref{sec:h1-c4c3}), where corroboration alone cannot resolve whether an
edge belongs in the final model.

\section{SysML v2 Generation}\label{sec:sysml-generation}

The final stage projects Stage~5's output, \texttt{elements.json} and
\texttt{sequence.json} (\cref{sec:agentic-synthesis}), into \sysml
textual notation, using the SysML v2 constructs described in
\cref{sec:sysml}. The \texttt{draw-sysml-structural-model} skill reads
\texttt{elements.json}: each unique component becomes a \texttt{part
def}, each interaction it sends or receives becomes a typed port on that
part, and each \texttt{(source, target, interaction)} triple becomes a
\texttt{connect} statement between the corresponding client and server
ports. The \texttt{draw-sysml-sequence-model} skill reads
\texttt{sequence.json}: each unique component becomes a lifeline
(\texttt{part def}), and each step becomes an ordered, labelled
\texttt{SynchronousCall} between the corresponding lifelines, with the
step number preserving the execution order decided during agentic
synthesis.

Component names are derived algorithmically from file paths, not from a
hardcoded mapping table, using a path-to-component resolution function
that handles \texttt{/src/}, \texttt{/include/}, \texttt{/extern/}, and
\texttt{/tests/} prefixes, CamelCase normalisation, and \texttt{Impl}
suffix stripping. This keeps the generator codebase-agnostic: the same
script produces meaningful lifeline names on CPSCore or any other source
tree following standard C++ project layout conventions.
